\documentclass[11pt]{amsart}
\usepackage{amsmath,amssymb,amsthm}
\usepackage{booktabs}
\usepackage[margin=1.2in]{geometry}
\usepackage[colorlinks=true,linkcolor=blue,citecolor=blue,urlcolor=blue]{hyperref}
\usepackage{microtype}
\emergencystretch=1.5em
\raggedbottom

\newtheorem{theorem}{Theorem}[section]
\newtheorem{proposition}[theorem]{Proposition}
\newtheorem{corollary}[theorem]{Corollary}
\newtheorem{lemma}[theorem]{Lemma}
\newtheorem{conjecture}[theorem]{Conjecture}
\newtheorem{question}[theorem]{Question}
\theoremstyle{definition}
\newtheorem{definition}[theorem]{Definition}
\newtheorem{example}[theorem]{Example}
\theoremstyle{remark}
\newtheorem{remark}[theorem]{Remark}
\newtheorem{observation}[theorem]{Observation}

\newcommand{\CC}{\mathbb{C}}
\newcommand{\ZZ}{\mathbb{Z}}
\newcommand{\QQ}{\mathbb{Q}}
\newcommand{\NN}{\mathbb{N}}
\newcommand{\Crit}{\operatorname{Crit}}
\newcommand{\Bal}{\operatorname{Balanced}}
\newcommand{\rank}{\operatorname{rank}}
\newcommand{\Mbar}{\overline{\mathcal{M}}}
\newcommand{\gfrak}{\mathfrak{g}}
\newcommand{\chat}{\hat{c}}
\DeclareMathOperator{\Aut}{Aut}

\hyphenation{de-scen-dent de-scen-dents wall-cross-ing sin-gu-lar-i-ty}

\title[The central-charge threshold for wall-crossing in open FJRW theory]{The
central-charge threshold for wall-crossing\\ in two-variable open FJRW theory}

\author{Bernd Johannes Wuebben}
\subjclass[2020]{14N35 (primary); 14J33, 53D45, 14B05 (secondary)}
\keywords{Open FJRW theory, Landau--Ginzburg models, wall-crossing,
central charge, simple singularities, boundary conditions}
\date{Initial version: August 24, 2026. This version: October 9, 2026}
\hypersetup{
  pdftitle={The central-charge threshold for wall-crossing in two-variable open FJRW theory},
  pdfauthor={Bernd Johannes Wuebben},
  pdfsubject={Open FJRW wall-crossing and the central-charge threshold},
  pdfkeywords={open FJRW theory, Landau--Ginzburg models, wall-crossing, central charge}
}

\begin{document}

\begin{abstract}
We prove that the genus-zero primary open FJRW invariants of Gross, Kelly and
Tessler for $x^r+y^s$ are independent of the choice of symmetric canonical
boundary conditions if and only if the central charge is less than one, that
is, exactly for the Fermat simple singularities $A$, $D_4$, $E_6$ and $E_8$.
The proof bounds their count of critical graphs, the graph types along which
the invariants jump, by the central charge, and passes from graphs to
invariants through the transitivity of their wall-crossing group.  That
transitivity needs one further generator, the vector field of the
anti-diagonal torus action, whenever a critical graph has one boundary
marking of each type; with it adjoined, the group acts freely and
transitively.  At and above the threshold, dependence first appears with two
interior insertions when both exponents are at least four, and with three
when the smaller is three; one insertion of descendent degree one produces it
for every pair.  At central charge one this vector field generates every
primary wall.  A B-model analogue, for open Saito potentials, is due to Maher.
\end{abstract}

\maketitle
\enlargethispage{3pt}

\section{Introduction}
\label{sec:intro}

\subsection*{Background}
For a nondegenerate quasi-homogeneous polynomial $W$ and a group $G$ of diagonal
symmetries, FJRW theory \cite{FJR13} assigns to the Landau--Ginzburg model
$(W,G)$ a cohomological field theory: a genus-zero $A$-model built from moduli
of $W$-spin curves and the Euler class of a Witten bundle. In their monograph
\cite{GKT} (whose numbering and page references we take from arXiv
version~3), Gross, Kelly and Tessler construct the \emph{open} analogue for the
two-variable Fermat model
\[
W=x^{r}+y^{s},\qquad G=\mu_{r}\times\mu_{s},\qquad r,s\geq 2,
\]
replacing closed $W$-spin curves by $W$-spin disks. The resulting open FJRW
invariants are integrals of multisections of a descendent Witten bundle over
moduli spaces that are real orbifolds with corners, and their central output is
an open mirror theorem: a generating function of the open invariants is a versal
deformation of $W$ whose oscillatory period integrals recover the closed
extended FJRW invariants, with the deformation parameters already flat
\cite[Theorem~0.3]{GKT}.

A structural feature distinguishes the open theory from its closed counterpart
and from the open $r$-spin theory of Buryak--Clader--Tessler
\cite{BCT21,BCT18}: the invariants are \emph{not} uniquely determined. Boundary conditions
for the Witten bundle over the corners of moduli space are constrained, but not
determined, by the positivity requirements of \cite{GKT}; the residual freedom
makes the open invariants jump as the boundary data cross certain walls. Gross,
Kelly and Tessler organize these jumps into a pro-nilpotent Landau--Ginzburg
wall-crossing group $G^{r,s}$ acting on the set of possible systems of
invariants, and state that this action is faithful and transitive
\cite[Theorems~0.7 and~4.27]{GKT}, so that the systems of open invariants form
a single torsor, while the period integrals, hence the closed theory, are
constant along it.  As explained below, transitivity requires one further
generator whenever a critical graph with one boundary marking of each type
occurs.

The source of every wall is a \emph{critical graph}: a stable $W$-spin
graph~$\Lambda$ whose Witten bundle has rank one above the dimension of its
moduli space, $\rank E_{\Lambda}(\mathbf d)=\dim\Mbar^{W}_{\Lambda}+1$
\cite[Definition~2.33]{GKT}. Such a graph indexes a one-parameter family of
codimension-one boundary strata through which zeros of the multisection can
escape. The wall-crossing Lie algebra $\gfrak^{r,s}$ of
\cite[Definition~4.23]{GKT} is a direct sum of one-dimensional pieces indexed
by critical graphs of nonzero bidegree (defined below); the excluded bidegree $(0,0)$ is
spanned by $x\partial_x-y\partial_y$, and we return to it below.

\subsection*{The prediction}
Whether critical graphs exist depends on $(r,s)$ and on whether one allows
descendents. In the descendent-free (\emph{primary}) theory,
\cite[Remark~0.9]{GKT} predicts that wall-crossing is absent exactly for the
simple singularities, and asserts this for an explicit list of pairs. As stated,
that remark contains a discrepancy between its list and its naming, which we
record and correct below. The prediction itself, that the primary open FJRW
invariants are independent of the boundary conditions if and only if $W$ is a
simple singularity, is left
open in \cite[Remark~0.9]{GKT}.

\subsection*{Results}
We prove the prediction as a threshold on the central charge.  Recall that the
\emph{central charge} of $W=x^{r}+y^{s}$, whose variables carry weights
$q_{x}=1/r$ and $q_{y}=1/s$, is
\begin{equation}
\label{eq:centralcharge}
\chat_{W}\;=\;(1-2q_{x})+(1-2q_{y})\;=\;2-\frac{2}{r}-\frac{2}{s}.
\end{equation}

Fix a finite set $I$ of labeled interior markings, each label $i$ carrying a
twist $(a_i,b_i)$ with $0\le a_i\le r-2$ and $0\le b_i\le s-2$.  A family of
GKT canonical boundary conditions is called \emph{symmetric} if it is invariant
under permutations of labels with the same twist.

\begin{theorem}[Central-charge criterion]
\label{thm:main}
Let $W=x^{r}+y^{s}$ with $r,s\geq2$ and central charge $\chat_{W}$ as
in~\eqref{eq:centralcharge}. The following are equivalent.
\begin{enumerate}
\item[(i)] $\chat_{W}<1$.
\item[(ii)] $\dfrac1r+\dfrac1s>\dfrac12$.
\item[(iii)] $W$ is a simple singularity: up to interchange of $x$ and $y$,
\[
(r,s)\in\{(n,2):n\geq2\}\cup\{(3,3),(3,4),(3,5)\},
\]
of type $A_{n-1}$, $D_{4}$, $E_{6}$, $E_{8}$ respectively.
\item[(iv)] For every finite labeled marking set $I$ as above, the genus-zero primary
open FJRW invariants of $(x^{r}+y^{s},\mu_{r}\times\mu_{s})$ are independent of
the choice of symmetric canonical boundary conditions.
\end{enumerate}
When these fail, that is when $\chat_{W}\geq1$, there is a finite marking set
$I$ for which the primary invariants depend on the boundary conditions.
\end{theorem}

The equivalence (i)$\Leftrightarrow$(ii)$\Leftrightarrow$(iii) is elementary and
classical \cite{Arnold72,Saito74}; the content is (iv). The threshold
$\chat_{W}=1$ is exactly the boundary between the simple singularities
($\chat_{W}<1$) and the parabolic (simple elliptic, $\chat_{W}=1$) and
hyperbolic ($\chat_{W}>1$) regimes; here \emph{hyperbolic} refers to the
triangle group $(2,r,s)$, not to Arnold's hyperbolic singularities
$T_{p,q,r}$, which are not quasi-homogeneous. Thus the presence of primary
wall-crossing detects, through open enumerative geometry, the classical ADE
threshold.

Beyond the dichotomy, the same count determines \emph{when} wall-crossing first
appears. For an interior insertion the relevant twists lie in the range
$0\leq a\leq r-2$, $0\leq b\leq s-2$; write $n$ for the number of interior
insertions.

\begin{theorem}[Least number of insertions]
\label{thm:firstwall}
Let $n_{*}(r,s)$ be the least number of primary interior insertions admitting a
critical graph. Then
\[
n_{*}(r,s)=
\begin{cases}
\infty, & \tfrac1r+\tfrac1s>\tfrac12\quad(\text{simple}),\\[2pt]
2, & \min(r,s)\geq 4,\\[2pt]
3, & \min(r,s)=3\ \text{and}\ \max(r,s)\geq 6.
\end{cases}
\]
The case $n_{*}=2$ is first realized by $W=x^{4}+y^{4}$ with two interior
insertions of twist $(2,2)$, the example of
\cite[Examples~4.9 and~5.5]{GKTsurvey}.
\end{theorem}

By the \emph{first wall} of $(r,s)$ we mean the set of critical graphs over
interior multisets of the minimal size $n_{*}(r,s)$.  By
Proposition~\ref{prop:reduction} below, $n_*(r,s)$ is also the least size of a
marking set over which the primary invariants depend on the boundary
conditions.

Finally, descendents change the picture entirely: they generate walls in every
case, including the simple singularities, as \cite[Remark~0.9]{GKT} anticipates.
The first such wall is universal.

\begin{proposition}[First descendent wall]
\label{prop:descendent}
For every $(r,s)$ with $r,s\geq2$, the data of a single interior insertion of
descendent degree one admit exactly one critical graph. Hence the open FJRW invariants with descendents
exhibit wall-crossing for every two-variable Fermat model, simple or not.
\end{proposition}

Thus the descendent-free theory is independent of the boundary conditions
exactly on the ADE locus, while with a single descendent the invariants
depend on them in every case.

\subsection*{The torus vector field}
The generator attached to a critical graph with $k_1+1$, respectively $k_2+1$,
boundary markings of the two types is a multiple of
$x^{k_1}y^{k_2}\bigl((k_2+1)x\partial_x-(k_1+1)y\partial_y\bigr)$ with a
coefficient in the deformation parameters; we call $(k_1,k_2)$ its
\emph{bidegree}.  GKT's Lie algebra $\gfrak^{r,s}$ contains the generators of
nonzero bidegree but excludes the bidegree $(0,0)$ \cite[Def.~4.23]{GKT},
whose generator is a multiple of the \emph{torus vector field}
$x\partial_x-y\partial_y$, the generator of the torus action
$(x,y)\mapsto(\lambda x,\lambda^{-1}y)$, which preserves $xy$ and
$dx\wedge dy$.  The critical graphs of bidegree $(0,0)$, those with one
boundary marking of each type, are exactly those of Euler weight zero
(Section~\ref{sec:count}); we call them \emph{weight-zero critical graphs}.
They occur for every $(r,s)$: a single interior insertion of twist $(0,0)$
with descendent degree one carries one.  At central charge one the first wall
consists of weight-zero critical graphs whose insertions carry the maximal
twist $(r-2,s-2)$, and $x^{r-2}y^{s-2}$ is the marginal deformation of $W$.
The description of the algebra in \cite[\S0.4 and Remark~4.24]{GKT} includes
the torus vector field, the jump argument of \cite[Theorem~5.9]{GKT} uses it,
and the realization argument of \cite[Theorem~5.13]{GKT} realizes the
corresponding motion of coefficients; the exclusion lies only in the explicit
definition.

\begin{proposition}[Enlargement of the Gross--Kelly--Tessler action]
\label{prop:enlargedaction}
Let $\widetilde G^{r,s}$ be the group generated by the wall-crossing group of
\cite{GKT} and the exponentials of the torus vector field multiplied by
rational multiples of the deformation monomials $u_{J,\mathbf d}$ whose data
$(J,\mathbf d)$ admit a weight-zero critical graph; equivalently
$\widetilde G^{r,s}=\exp(\widetilde\gfrak^{\,r,s})$ of
\eqref{eq:enlargedgroup}.  It preserves the set of systems of symmetric
canonical open FJRW invariants with marking set $I$ and acts on it freely and
transitively.
\end{proposition}

The wall-crossing group of \cite{GKT} alone does not act transitively as soon
as $I$ carries a weight-zero critical graph; the smallest case is a single
insertion of twist $(0,0)$ with descendent degree one
(Example~\ref{ex:failure}).  The proof of preservation also extends the
period argument of \cite{GKT} to all residues (Remark~\ref{rem:marginalGKT}).
In the parabolic regime the torus vector field generates every primary wall:
every primary critical graph has weight zero, because equality in the bound
of the Method paragraph below forces maximal twists and vanishing residues
(Proposition~\ref{prop:parabolic}).

\subsection*{The wall generators}
Gross, Kelly and Tessler write the wall-crossing Lie algebra in terms of the
vector fields above \cite[eqn.~(4.19)]{GKT} and observe that they preserve
$dx\wedge dy$ \cite[p.~89]{GKT}; conversely, every derivation preserving
$dx\wedge dy$ and the ideal $(xy)$ is a combination of these fields
\cite[Remark~4.24]{GKT}, and Maher's algebra \cite[(5.2.1) and
Proposition~5.2.1]{Maher} consists of the same fields.  Each of them is the
Hamiltonian vector field of the monomial $x^{k_1+1}y^{k_2+1}$ for
$dx\wedge dy$, and the structure constants of the bracket are the
determinants of the numbers of boundary markings, as in the tropical vertex
Lie algebra of the logarithmic form $d\log x\wedge d\log y$; the two Lie
algebras are
nevertheless not isomorphic as graded Lie algebras (Section~\ref{sec:hamiltonian}).

Together, Theorems~\ref{thm:main} and~\ref{thm:firstwall} and
Propositions~\ref{prop:descendent} and~\ref{prop:enlargedaction} describe when
walls exist and how the wall-crossing group acts.  The slope-ordered
factorization of the wall-crossing group, with rays indexed by the numbers of
boundary markings $(k_1+1,k_2+1)$ of critical graphs, is constructed in the
companion paper \cite{companion}.

\subsection*{Method}
Nothing here rebuilds the machinery of \cite{GKT}.  GKT's count
\cite[Notation~2.40]{GKT} says that a critical graph with interior data
$(J,\mathbf d)$ exists if and only if
$N(J,\mathbf d)=\ell_1+\ell_2-|J|+1+\sum_i d_i\ge1$, where $\ell_1$ and $\ell_2$
are the carries in the twist sums $\sum a_i$ and $\sum b_i$.  Since
Neveu--Schwarz twists are at most $(r-2,s-2)$, the single inequality of
Lemma~\ref{lem:bound} gives $\ell_1+\ell_2\le |J|\,\chat_W$, so
$N(J,\mathbf 0)\le 0$ when $\chat_W<1$.  When $\chat_W\ge1$, $n$ copies of the
maximal twist give $N=1$ for suitable $n$; since $N$ is nondecreasing in each
twist, the least such $n$ is $n_*$.  A single insertion of descendent degree
$d$ has $N=d$, which gives Proposition~\ref{prop:descendent}.  The step from
graphs to invariants is the reduction to critical graphs
(Proposition~\ref{prop:reduction}): the generator of a primary critical graph
changes a primary coefficient of the potential; if there is none, every
generator carries a descendent variable, and transitivity forces all systems
to agree.  For Proposition~\ref{prop:enlargedaction}, preservation holds
because every generator is divergence-free and tangent to the axes, so it
preserves the integral of $e^{W/\hbar}$ over the positive quadrant, whose
coefficients are, up to fixed nonzero factors, GKT's polynomials
$\mathcal A(J,\mathbf d,\nu)$; this replaces the period formula of
\cite[Theorem~4.9]{GKT}, which is stated only for residues at most $r-2$ and
$s-2$.  Freeness follows by looking at a divisibility-minimal deformation
monomial in the support of a nonzero element $v$ of the enlarged Lie
algebra, for $g=\exp v$.  Transitivity is proved by induction
over divisibility: at each monomial the difference of two systems lies in the
kernel of one nonzero linear functional on the coefficients of the balanced
graphs, and the vectors by which the critical graphs move these coefficients
span that kernel.  All numerical claims are proved in
\S\S\ref{sec:proof}--\ref{sec:parabolic}; finite exact checks are described in
\S\ref{sec:cert}.

\subsection*{Relation to prior work}
The uniqueness of the primitive form for simple singularities is classical
\cite{Saito83,NoumiYamada98}, and is the closed-mirror shadow of the statement
that these are the potentials without wall-crossing. On the open side,
Basalaev--Buryak constructed open Saito potentials for the $A$- and $D$-series
by solving the open WDVV equations \cite{BB21}, and Maher computed open
Saito potentials for the Fermat form $x^3+y^3$ of $D_4$ and for the
chain-type $E_7$ \cite[Examples~4.2.3 and~4.2.4]{Maher}; the present criterion
explains, uniformly, why the primary open FJRW theory is unambiguous on the
whole ADE locus. The wall-crossing group of \cite{GKT} is drawn, ``in spirit,'' by analogy
with the tropical vertex group of Kontsevich--Soibelman \cite{KS06} and its role
in Fano mirror symmetry \cite{GPS10,Gross10} (\cite[Remark~0.8]{GKT}); separately, a
tropical version of open FJRW theory is raised as an open problem in
\cite[\S6]{GKTsurvey}, there through tropical moduli. The classification of critical
graphs organized here is the starting point for the slope-ordered
factorization of the wall-crossing group in \cite{companion}.

Maher proves a B-model analogue of the threshold: for invertible polynomials,
an axiomatically defined primary wall-crossing group of open Saito potentials
is trivial exactly for the simple singularities
\cite[Proposition~1.0.5 and Corollary~5.2.6]{Maher}.  His wall-crossing
algebra \cite[(5.2.1)]{Maher} includes the torus vector field, and for
simple elliptic singularities of rank two or three he shows that the primary
wall-crossing elements are exponentials of combinations of torus fields
(in rank two, of $x\partial_x-y\partial_y$) with coefficients powers of the
marginal parameter satisfying an integrality condition
\cite[Propositions~1.0.5 and~5.2.10]{Maher}; Proposition~\ref{prop:parabolic}
below is the corresponding statement in the open FJRW theory of \cite{GKT}.
The A-model statement (i)$\Leftrightarrow$(iv) of Theorem~\ref{thm:main}, the
first-wall count of Theorem~\ref{thm:firstwall} and the descendent wall are,
to our knowledge, not in \cite{Maher}; we give an elementary proof in the
Fermat case.

\section{The count of critical graphs}
\label{sec:count}

We recall only what is needed and refer to \cite{GKT} for all definitions. Fix
$W=x^{r}+y^{s}$. A genus-zero smooth $(r,s)$-disk carries interior marked points
with twists $(a_{i},b_{i})$, $0\leq a_{i}\leq r-1$, $0\leq b_{i}\leq s-1$, and a
descendent vector $\mathbf d=(d_{i})\in\ZZ_{\geq0}^{J}$. The interior insertions
contributing to the primary potential lie in the Neveu--Schwarz range
$0\leq a_{i}\leq r-2$, $0\leq b_{i}\leq s-2$ \cite[eqn.~(0.5)]{GKT}.

Given a set $J$ of interior markings with twists $(a_{i},b_{i})_{i\in J}$ and
descendents $\mathbf d=(d_{i})_{i\in J}$, define, following
\cite[Notation~2.30]{GKT},
\begin{equation}
\label{eq:ells}
\sum_{i\in J}a_{i}=r(J)+\ell_{1}r,\qquad
\sum_{i\in J}b_{i}=s(J)+\ell_{2}s,
\end{equation}
where $r(J)\in\{0,\dots,r-1\}$ and $s(J)\in\{0,\dots,s-1\}$ are the residues, so
that
\begin{equation}
\label{eq:elldef}
\ell_{1}=\Big\lfloor\tfrac{1}{r}\textstyle\sum_{i\in J}a_{i}\Big\rfloor,\qquad
\ell_{2}=\Big\lfloor\tfrac{1}{s}\textstyle\sum_{i\in J}b_{i}\Big\rfloor.
\end{equation}
The classification of balanced and critical graphs with interior markings
exactly $J$ is \cite[Notation~2.40]{GKT}: setting
\begin{equation}
\label{eq:N}
N(J,\mathbf d)\;=\;\ell_{1}+\ell_{2}-|J|+1+\sum_{i\in J}d_{i},
\end{equation}
one has $\bigl|\Bal(J,\mathbf d)\bigr|=\max\bigl(N(J,\mathbf d)+1,\,0\bigr)$ and
\begin{equation}
\label{eq:critcount}
\bigl|\Crit(J,\mathbf d)\bigr|=\max\bigl(N(J,\mathbf d),\,0\bigr).
\end{equation}
Here the balanced graphs are the disk types whose invariants occur as
coefficients of the open potential; the critical graphs index the possible
codimension-one jumps between those coefficients.
In particular a critical graph with interior markings $J$ and descendent vector
$\mathbf d$ exists if and only if $N(J,\mathbf d)\geq1$. (Throughout,
$J\neq\emptyset$. For $J=\emptyset$, the balanced equation has precisely the
solutions $(k_{1},k_{2})=(r,0),(0,s)$, the two boundary-only graphs appearing
in \cite[Example~3.4]{GKT}; the critical equation forces $(0,0)$, whose two
boundary markings are unstable. Thus \eqref{eq:critcount} is not asserted in
the boundary-only case.)

\subsection*{Weights and the enlarged algebra}
For a nonempty subset $J\subseteq I$ and descendent degrees
$\mathbf d=(d_i)_{i\in J}$, set
\begin{gather}
 r(J)\equiv\sum_{i\in J}a_i\pmod r,\quad 0\le r(J)<r,
 \qquad
 s(J)\equiv\sum_{i\in J}b_i\pmod s,\quad 0\le s(J)<s,\label{eq:residues}\\
 m(J,\mathbf d)
 =rs+\sum_{i\in J}\bigl(sa_i+rb_i+rs(d_i-1)\bigr),
 \qquad w(J,\mathbf d)=m(J,\mathbf d)-rs.\label{eq:massweight}
\end{gather}
Here $m(J,\mathbf d)$ is the quantity of \cite[Notation~2.30]{GKT}.  The generator attached to a critical graph with
$k_1+1$, respectively $k_2+1$, boundary markings of the two types is a
deformation-monomial multiple of
$x^{k_1}y^{k_2}\bigl((k_2+1)x\partial_x-(k_1+1)y\partial_y\bigr)$, whose weight
under the Euler vector field $sx\partial_x+ry\partial_y$ of $W$ is
$sk_1+rk_2$, and the critical condition \cite[Proposition~2.39]{GKT} reads
\begin{equation}
 sk_1+rk_2=w(J,\mathbf d).
 \label{eq:criticalweight-intro}
\end{equation}
We therefore call $w$ the \emph{Euler weight}.  Write $t_{a,b,d}$ for the
deformation variable associated with an insertion $\tau_d^{(a,b)}$ and
$u_{J,\mathbf d}=\prod_{i\in J}t_{a_i,b_i,d_i}$.  The coefficient ring
$A_{I,\mathrm{sym}}$ is the symmetric GKT deformation ring in which the total
multiplicity of each twist is bounded by the number of labels of that twist in
$I$; in particular every positive-degree element is nilpotent.
For a balanced graph, $k_1$ and $k_2$ denote its numbers of boundary markings
of the two types; for a critical graph with $k_1+1$ and $k_2+1$ boundary
markings, $(k_1,k_2)$ is its bidegree, as in the introduction.  Since
$k_1,k_2\ge0$, \eqref{eq:criticalweight-intro} shows that a critical graph has
Euler weight zero exactly when its bidegree is $(0,0)$, that is, when it is a
weight-zero critical graph in the sense of the introduction.

The \emph{enlarged} algebra and group of Proposition~\ref{prop:enlargedaction} over
$A_{I,\mathrm{sym}}$ are
\begin{equation}
\begin{aligned}
 \widetilde{\gfrak}^{\,r,s}_{A_{I,\mathrm{sym}}}
 &=\gfrak^{r,s}_{A_{I,\mathrm{sym}}}
 \oplus\!\!\sum_{\substack{\varnothing\ne J\subseteq I,\;\mathbf d\ge0\\
 w(J,\mathbf d)=0,\;r(J)=s(J)=0}}
 \QQ\,u_{J,\mathbf d}(x\partial_x-y\partial_y),\\
 \widetilde G^{r,s}_{A_{I,\mathrm{sym}}}
 &=\exp\bigl(\widetilde{\gfrak}^{\,r,s}_{A_{I,\mathrm{sym}}}\bigr).
\end{aligned}
\label{eq:enlargedgroup}
\end{equation}
The exponential is finite over $A_{I,\mathrm{sym}}$.  This notation keeps
GKT's original positive-weight group distinct from the enlargement used here.






\section{The torus vector field and the enlarged action}
\label{sec:enlarged}

We recall the objects of \cite{GKT} on which the enlarged action acts.  A
system of open invariants bounded by $I$ is a vector
$\nu=(\nu_{\Gamma,\mathbf d})$ indexed by the balanced graphs $\Gamma$ with
internal labels in $I$ and their descendent vectors $\mathbf d$
\cite[(0.3)]{GKT}.  For $J\subseteq I$ and $\mathbf d$, the chamber-index
polynomial of \cite[Notation~3.42]{GKT}, which we call the chamber-index
polynomial, is
\begin{equation}
\label{eq:Apoly}
\mathcal A(J,\mathbf d,\nu)=\sum_{h=1}^{|J|}\frac1{h!}
\sum_{J_1\sqcup\cdots\sqcup J_h=J}\ \sum_{\Theta_1,\dots,\Theta_h}
\frac{\Gamma\bigl(\frac{1+\sum_i k_1(\Theta_i)}{r}\bigr)
      \Gamma\bigl(\frac{1+\sum_i k_2(\Theta_i)}{s}\bigr)}
     {\Gamma\bigl(\frac{1+r(J)}{r}\bigr)\Gamma\bigl(\frac{1+s(J)}{s}\bigr)}
\prod_{i=1}^h\nu_{\Theta_i,\mathbf d|_{J_i}},
\end{equation}
the sum over ordered decompositions of $J$ into nonempty blocks and over
balanced graphs $\Theta_i$ with internal labels $J_i$, where $k_1(\Theta_i)$
and $k_2(\Theta_i)$ are the numbers of boundary markings of the two types and
$\Gamma$ is the Gamma function.  A
\emph{chamber index} is a system with $\nu_{\Theta,\mathbf 0}=1$ when the
balanced graph $\Theta$ has one internal marking, $\nu_{\Theta,\mathbf 0}=-1$
when it has none,
$\mathcal A(J,\mathbf d,\nu)=(-1)^d$ when $J=\{j\}$ and $\mathbf d=(d)$, and
$\mathcal A(J,\mathbf d,\nu)=0$ when $|J|\ge2$ and
$m(J,\mathbf d)\ge s\,r(J)+r\,s(J)$ \cite[Definition~3.47]{GKT}.  By
\cite[Corollary~5.14]{GKT}, the chamber indices are exactly the systems
realized by canonical boundary conditions; we call them \emph{realizable}.

\begin{remark}[The torus vector field in {\cite{GKT}}]
\label{rem:marginalGKT}
Besides the descendent singleton of twist $(0,0)$, examples of weight-zero
critical graphs are \cite[Example~2.38]{GKT} and
\cite[Example~4.9]{GKTsurvey}.  With the Lie algebra of
\cite[Definition~4.23]{GKT}, which omits the bidegree $(0,0)$, the
transitivity in \cite[Theorems~0.7 and~4.27]{GKT} fails as soon as $I$
carries such a graph: by the freeness in
Proposition~\ref{prop:enlargedaction}, the only element of
$\widetilde G^{r,s}$ carrying a realizable system $\nu$ to
$\exp\bigl(c\,u(x\partial_x-y\partial_y)\bigr)\nu$, with $c\ne0$ and $u$ the
deformation monomial of the graph, is that element itself, which does not
lie in $G^{r,s}$; Example~\ref{ex:failure} is the smallest case.  On the
family of balanced graphs of such a $u$, the element written in
\cite[proof of Theorem~4.27, Step~3]{GKT} is a multiple of the torus vector
field.  The jump automorphisms in the proof of \cite[Theorem~5.9]{GKT}, over
the labeled ring of \cite[Definition~4.6]{GKT}, lie in the analogous
enlargement of GKT's group over that ring, because the jump sets
$Q_0(A,\mathbf d)$ there impose only the congruences of
\cite[Definition~4.23]{GKT} and \eqref{eq:criticalweight-intro}, which admit
$(k_1,k_2)=(0,0)$ when $r(A)=s(A)=0$.  The equality of the chamber indices with the
realizable systems \cite[Corollary~5.14]{GKT}, proved through
\cite[Theorem~5.13]{GKT}, does not depend on transitivity: at
\cite[(5.50)]{GKT} the two required values of $\mathcal A(J,\mathbf d,-)$ agree
because both are prescribed by \cite[Definition~3.47]{GKT}.  The preservation
step of \cite[proof of Theorem~4.27, Step~1]{GKT} uses the period formula of
\cite[Theorem~4.9]{GKT}, which is stated for residues at most $r-2$ and
$s-2$; the Laplace integral in the proof below covers the residues $r-1$ and
$s-1$ as well.
\end{remark}

\begin{example}[Failure of transitivity without the torus vector field]
\label{ex:failure}
Let $I=\{1\}$ consist of one label of twist $(0,0)$ and consider descendent
degree one, so $u=t_{0,0,1}$, $J=\{1\}$, $r(J)=s(J)=0$ and $N(J,(1))=1$.  The
balanced graphs have boundary markings $(0,s)$ and $(r,0)$, and the unique
critical graph has bidegree $(0,0)$.  By \eqref{eq:Apoly}, the chamber-index
condition $\mathcal A(J,(1),\nu)=-1$ reads
\[
 \frac{\Gamma(\frac1r)\Gamma(\frac{1+s}s)}{\Gamma(\frac1r)\Gamma(\frac1s)}\,
 \nu_{(0,s)}
 +\frac{\Gamma(\frac{1+r}r)\Gamma(\frac1s)}{\Gamma(\frac1r)\Gamma(\frac1s)}\,
 \nu_{(r,0)}
 =\frac{\nu_{(0,s)}}{s}+\frac{\nu_{(r,0)}}{r}=-1 ,
\]
a single linear condition on two coefficients.  Every solution extends to a
chamber index, since for $|I|=1$ the conditions for different descendent
degrees involve disjoint coordinates, and is therefore realizable
\cite[Corollary~5.14]{GKT}.  A generator of $\gfrak^{r,s}$ changes
the $u$-coefficients of a potential only if its coefficient divides $u$, that
is, equals $u$; the generators with coefficient $u$ are attached to the
critical graphs of $(J,(1))$, and the only one has bidegree $(0,0)$, which
\cite[Definition~4.23]{GKT} excludes.  Hence no element of $G^{r,s}$ relates
two realizable systems with different values of $\nu_{(r,0)}$.  The element
$\exp\bigl(c\,u(x\partial_x-y\partial_y)\bigr)$ of the enlarged group moves the
pair $(\nu_{(0,s)},\nu_{(r,0)})$ along the solution line, by a multiple of
$(-s,r)$.
\end{example}

\begin{proof}[Proof of Proposition~\ref{prop:enlargedaction}]
\emph{Closure and finiteness.}  The enlarged algebra is closed under the
bracket: for a weight-zero coefficient $u$ and a generator $u'X_{k_1,k_2}$,
$[u(x\partial_x-y\partial_y),u'X_{k_1,k_2}]=(k_1-k_2)uu'X_{k_1,k_2}$, and
$uu'$ satisfies the residue and weight conditions of $u'$ because those of
$u$ vanish; brackets of positive-weight generators are as in
\cite[Prop.~4.25]{GKT}.  Every generator has a coefficient in the ideal
generated by the deformation variables, which is nilpotent in
$A_{I,\mathrm{sym}}$, so the exponential is finite.  The argument of Step~1
in the proof of \cite[Thm.~4.27]{GKT} applies verbatim to the weight-zero
generators and shows that $e^vW^\nu$ is again the potential of a symmetric
system of balanced coefficients.

\emph{Preservation.}  Take $\hbar<0$ and expand, coefficientwise, the
positive-quadrant Laplace integral
\[
 \mathcal L_\hbar(W)=\int_0^\infty\!\int_0^\infty e^{W/\hbar}\,dx\,dy.
\]
Its monomial moments are
\[
 \int_0^\infty\!\!\int_0^\infty x^{\alpha}y^{\beta}
 e^{(x^r+y^s)/\hbar}\,dx\,dy
 =\frac1{rs}(-\hbar)^{\frac{\alpha+1}r+\frac{\beta+1}s}
 \Gamma\Bigl(\frac{\alpha+1}r\Bigr)\Gamma\Bigl(\frac{\beta+1}s\Bigr).
\]
Write $W=x^r+y^s+V$ and expand $e^{V/\hbar}=\sum_h V^h/(h!\,\hbar^h)$.  The
$u_{J,\mathbf d}$-coefficient of $V^h$ is a sum over ordered decompositions
$J=J_1\sqcup\cdots\sqcup J_h$ of products of coefficients of $V$, each
attached to a balanced graph with internal labels $J_i$.  Since
$\sum_i\bigl(s\,k_1(\Theta_i)+r\,k_2(\Theta_i)\bigr)=m(J,\mathbf d)+(h-1)rs$, all
these terms carry the same power of $-\hbar$; the sign $(-1)^h$ from
$\hbar^{-h}$ and the signs $(-1)^{|J_i|-1}$ of the coefficients of the
potential combine to $(-1)^{|J|}$, and the moment of the product monomial is
the Gamma product in \eqref{eq:Apoly}.  Hence the $u_{J,\mathbf d}$-coefficient
of $\mathcal L_\hbar(W^\nu)$ is a nonzero multiple, depending only on
$(J,\mathbf d)$, of $\mathcal A(J,\mathbf d,\nu)$, for every pair of residues
$r(J),s(J)$.  This is the termwise calculation in the proof of
\cite[Theorem~4.9]{GKT}, with the positive quadrant in place of the
good-basis cycles, carried out over the labeled ring of
\cite[Definition~4.6]{GKT}; the injective map $\psi_I$ of
\cite[Lemma~4.14]{GKT} sends $W^{\nu,\mathrm{sym}}$ to $W^\nu$
\cite[Corollary~4.15]{GKT} and is compatible with the action, so over
$A_{I,\mathrm{sym}}$ the coefficient acquires the factor $1/|\Aut(A)|$ of
\cite[(4.11)]{GKT}.  Every enlarged generator
is divergence-free (Lemma~\ref{lem:hamvf}) and tangent to both coordinate
axes, since in \eqref{eq:Xab} the coefficient of $\partial_x$ vanishes on
$x=0$ and that of $\partial_y$ on $y=0$.  Integration by parts therefore gives
$\mathcal L_\hbar(e^vW)=\mathcal L_\hbar(W)$.  Thus all the
$\mathcal A$-conditions are preserved, including the residue classes $r-1$
and $s-1$, which the good-basis periods of \cite[\S4.1]{GKT} do not detect.
The normalizations of Definition~3.47 are also preserved: they prescribe the
coefficients of balanced graphs with at most one internal marking and no
descendents, and a generator changes the coefficient of a deformation
monomial only if the deformation monomial of some critical graph divides
it; a
primary critical graph has at least two interior markings, and a critical
graph with one interior marking carries a descendent.  Since the chamber indices are exactly the realizable
systems, this proves preservation.

\emph{The wall vectors.}  Fix a deformation monomial $u=u_{J,\mathbf d}$ and
write $N=N(J,\mathbf d)$.  The balanced disk types over $(J,\mathbf d)$ are
indexed by $0\le p\le N$ and have boundary degrees
\[
 (r(J)+pr,\ s(J)+(N-p)s).
\]
Let $V_u=\bigoplus_{p=0}^{N}\QQ e_p$ be their coefficient space, where $e_p$
changes only the coefficient of the $p$th balanced type.  For $1\le p\le N$,
the critical graph between the $(p-1)$st and $p$th balanced types acts on
$V_u$, through $x^r+y^s$, by
\[
 v_p=\alpha_pe_p-\beta_pe_{p-1},\qquad
 \alpha_p=r\bigl(s(J)+(N-p)s+1\bigr),\qquad
 \beta_p=s\bigl(r(J)+(p-1)r+1\bigr).
\]
Both coefficients are positive, and reading successively the
$e_N,e_{N-1},\ldots,e_1$ coordinates shows that the $v_p$ are independent.  If
$u$ carries a weight-zero critical graph, the weight equation forces $N=1$ and
$r(J)=s(J)=0$; its only wall vector is $v_1=re_1-se_0$, the vector of the
torus vector field adjoined in \eqref{eq:enlargedgroup}.

\emph{Freeness.}  Let $g=\exp(v)$ with $v\ne0$, and let $u$ be minimal for
divisibility among the deformation monomials in the support of $v$.  The
change of the $u$-coefficients of any potential under $g$ is $v_u(x^r+y^s)$,
where $v_u$ is the $u$-component of $v$: every other contribution carries a
product of a support monomial with a further deformation monomial, which by
minimality is not $u$.  The summands of $v_u$ act by the independent vectors
$v_p$, so $v_u(x^r+y^s)\ne0$ and $g$ moves every potential.  The same
argument applies modulo the monomials outside any finite divisor-closed set
of deformation monomials.

\emph{Transitivity.}  Let $S$ be a finite set of deformation monomials that
contains every divisor of each of its elements, and pass to the quotient that
retains exactly the coefficients indexed by $S$.  Order $S$ by a linear
extension of proper divisibility, and suppose two realizable systems agree on
every proper divisor of $u\in S$.  They satisfy the same prescribed value of
$\mathcal A(J,\mathbf d,-)$: by Definition~3.47(2) of \cite{GKT} if $|J|=1$,
and by its condition~(3) if $|J|\ge2$, which applies because the existence of
a balanced type forces $m(J,\mathbf d)\ge s\,r(J)+r\,s(J)$.  Every term of
\eqref{eq:Apoly} with two or more blocks already agrees, so their difference
in $V_u$ lies in the kernel of the remaining one-block term, a nonzero linear
functional $L_u:V_u\to\QQ$ whose coefficients are the Gamma factors computed
above; thus $\dim\ker L_u=N$.  The Laplace invariance gives $L_u(v_p)=0$, so
the $v_p$ form a basis of $\ker L_u$, and a unique combination of the
critical-graph generators with coefficient $u$ removes the difference in
$V_u$.  Acting by that combination can alter only coefficients indexed by
strict multiples of $u$, so it does not disturb earlier steps.  Induction
proves transitivity in the quotient defined by $S$.  By freeness, the element
so constructed is the unique one in the quotient group carrying the first
system to the second, so the elements obtained for $S\subset S'$ are
compatible.  For each bound on the degree $k_1+k_2$ of the generators $X_{k_1,k_2}$
only finitely many deformation monomials occur, because the weight equation
\eqref{eq:criticalweight-intro} bounds the descendent degrees.  The compatible
elements therefore define an element of the group \eqref{eq:enlargedgroup},
which, as in \cite[Def.~4.19]{GKT}, is the inverse limit of its truncations by
the degree $k_1+k_2$; it carries one system to the other, which proves
transitivity over $A_{I,\mathrm{sym}}$.
\end{proof}



\section{Proof of the central-charge criterion}
\label{sec:proof}

Both directions of Theorem~\ref{thm:main} rest on the following reduction from
invariants to graph counts.

\begin{proposition}[Reduction to critical graphs]
\label{prop:reduction}
Fix a finite marking set $I$ and work over the coefficient ring
$A_{I,\mathrm{sym}}$ of \S\ref{sec:count}.
\begin{enumerate}
\item[(a)] If no critical graph with descendent vector $\mathbf0$ has interior
markings contained in $I$, then
$\widetilde G^{r,s}_{A_{I,\mathrm{sym}}}$ fixes every primary open FJRW
invariant; the primary invariants are independent of the boundary
conditions.
\item[(b)] If some critical graph with descendent vector $\mathbf d$ has
interior markings $J\subseteq I$, then there are two systems of symmetric
canonical open FJRW invariants over $I$ that differ in an invariant with interior markings $J$ and
descendent vector $\mathbf d$; if $\mathbf d=\mathbf 0$, this is a primary
invariant.
\end{enumerate}
\end{proposition}

\begin{proof}
By \cite[Definition~4.23]{GKT}, the original wall-crossing Lie algebra
$\gfrak^{r,s}_{A_{I,\mathrm{sym}}}$ is spanned by the summands
corresponding to critical graphs with $(k_{1},k_{2})\ne(0,0)$.
Equation~\eqref{eq:enlargedgroup} adds the omitted $(0,0)$ summands. After passage
to symmetric coefficients, critical graphs with the same twist--descendent
multiset determine the same summand. For a critical graph with interior
markings $J$, descendents $\mathbf d$, and bidegree
$(k_{1},k_{2})$, its generator is
\begin{equation}
\label{eq:generator}
v_{\Lambda}=u_{J,\mathbf d}\,x^{k_{1}}y^{k_{2}}
\bigl((k_{2}+1)x\partial_{x}-(k_{1}+1)y\partial_{y}\bigr),\qquad
u_{J,\mathbf d}=\prod_{i\in J}t_{a_{i},b_{i},d_{i}} .
\end{equation}
Here $(k_{1},k_{2})$ is the bidegree used in
\cite[Definition~4.23]{GKT}: the critical graph has $k_{1}{+}1$ and
$k_{2}{+}1$ boundary marked points of the two types
\cite[eqn.~(2.31), Proposition~2.39]{GKT}, and
the bidegree satisfies the critical condition
$sk_{1}+rk_{2}=w(J,\mathbf d)$ of \eqref{eq:criticalweight-intro}. Positive
weight gives a summand of $\gfrak^{r,s}$, while $(k_{1},k_{2})=(0,0)$ gives the
adjoined torus vector field. GKT's group
$G^{r,s}_{A_{I,\mathrm{sym}}}=\exp(\gfrak^{r,s}_{A_{I,\mathrm{sym}}})$ acts on
the symmetric chamber indices of \cite[Definition~3.47]{GKT} by
$g(W^{\nu})=W^{g(\nu)}$ \cite[Theorem~4.27]{GKT} (the action itself is
reproved in Proposition~\ref{prop:enlargedaction});
Proposition~\ref{prop:enlargedaction} extends the action transitively to
$\widetilde G^{r,s}_{A_{I,\mathrm{sym}}}$.

For (a): if there is no $\mathbf0$-critical graph, every generator
\eqref{eq:generator} has some $d_{i}>0$, hence $u_{J,\mathbf d}$ is divisible by
a variable $t_{a,b,d}$ with $d>0$. Acting on a potential
$W^{\nu}=x^{r}+y^{s}+\sum_{0\le a\le r-2,\ 0\le b\le s-2}t_{a,b,0}\,x^{a}y^{b}+\cdots$
\cite[eqn.~(0.5)]{GKT}, such a generator raises the total degree in the
positive-descendent variables and therefore leaves unchanged the coefficients of
the purely primary monomials $u_{J',\mathbf 0}\,x^{\alpha}y^{\beta}$. Since
each primary open invariant $\nu_{\Gamma,\mathbf 0}$, the coefficient attached
to a graph $\Gamma$ in \cite[eqns.~(0.4)--(0.5)]{GKT}, is, up to the fixed
nonzero sign and automorphism factor in the symmetric potential, such a
coefficient, it is fixed by every generator, hence by
$\widetilde G^{r,s}_{A_{I,\mathrm{sym}}}$. Transitivity of the enlarged action
then makes these primary invariants equal in every system.

For (b): a critical graph $\Lambda$ with data $(J,\mathbf d)$ furnishes a
nonzero generator $v_{\Lambda}$ with coefficient $u_{J,\mathbf d}$, a monomial
in primary variables only when $\mathbf d=\mathbf 0$. Fix any system
$\nu$ of symmetric canonical invariants, which exists by
\cite[Theorem~0.1]{GKT}. The coefficient of the deformation monomial
$u_{J,\mathbf d}$ in
$\exp(v_{\Lambda})(W^{\nu})-W^{\nu}$ is exactly
$v_{\Lambda}(x^{r}+y^{s})$, namely
\[
u_{J,\mathbf d}\,\bigl(r(k_{2}+1)\,x^{k_{1}+r}y^{k_{2}}
-s(k_{1}+1)\,x^{k_{1}}y^{k_{2}+s}\bigr),
\]
which is nonzero. Proposition~\ref{prop:enlargedaction} ensures that
$\exp(v_{\Lambda})\cdot\nu$ is another system of symmetric canonical invariants
over $I$. The two systems therefore differ in an invariant with interior
markings $J$ and descendent vector $\mathbf d$.
\end{proof}


The equivalences (i)$\Leftrightarrow$(ii)$\Leftrightarrow$(iii) are immediate:
\eqref{eq:centralcharge} gives $\chat_{W}<1\Leftrightarrow\frac2r+\frac2s>1$, and
$\frac1r+\frac1s>\frac12$ characterizes the two-variable Fermat simple
singularities, whose only types are $A_{n-1}=(n,2)$, $D_{4}=(3,3)$,
$E_{6}=(3,4)$ and $E_{8}=(3,5)$ (there is no Fermat $E_{7}$, as $E_{7}=x^{3}+xy^{3}$
is not of the form $x^{r}+y^{s}$; likewise $D_{n}=x^{n-1}+xy^{2}$ is Fermat only
for $n=4$, where $x^{3}+y^{3}\cong D_{4}$). It remains to prove
(i)$\Leftrightarrow$(iv), for which, by Proposition~\ref{prop:reduction}, it
suffices to decide the existence of a primary ($\mathbf d=\mathbf0$) critical
graph. The whole matter rests on one inequality.

\begin{lemma}
\label{lem:bound}
Let $J$ be a set of $n=|J|$ interior markings with Neveu--Schwarz twists,
$0\leq a_{i}\leq r-2$ and $0\leq b_{i}\leq s-2$. Then, with $\ell_{1},\ell_{2}$
as in \eqref{eq:elldef},
\[
\ell_{1}+\ell_{2}\;\leq\;n\,\chat_{W}.
\]
Consequently the primary count \eqref{eq:N} obeys
$N(J,\mathbf 0)\leq 1-n(1-\chat_{W})$.
\end{lemma}

\begin{proof}
Dropping the floor functions and using the twist bounds,
\[
\ell_{1}+\ell_{2}
=\Big\lfloor\tfrac1r\textstyle\sum a_{i}\Big\rfloor
+\Big\lfloor\tfrac1s\textstyle\sum b_{i}\Big\rfloor
\leq\frac{\sum a_{i}}{r}+\frac{\sum b_{i}}{s}
\leq\frac{n(r-2)}{r}+\frac{n(s-2)}{s}
=n\Big(2-\tfrac2r-\tfrac2s\Big)=n\,\chat_{W}.
\]
The second claim is $N(J,\mathbf 0)=\ell_{1}+\ell_{2}-n+1\leq n\chat_{W}-n+1$.
\end{proof}

\begin{proof}[Proof of (i)$\Rightarrow$(iv)]
Assume $\chat_{W}<1$. For any $J$ with $n=|J|\geq1$ interior markings,
Lemma~\ref{lem:bound} gives $\ell_{1}+\ell_{2}\leq n\chat_{W}<n$; as
$\ell_{1}+\ell_{2}$ is an integer strictly below the integer $n$, we get
$\ell_{1}+\ell_{2}\leq n-1$, i.e.\ $N(J,\mathbf 0)\leq0$. By
\eqref{eq:critcount} there is no primary critical graph, for any $J$ and any
finite $I$. Proposition~\ref{prop:reduction}(a) then says the primary invariants
are independent of the boundary conditions for every $I$. (The case $n=0$ is vacuous: there is no interior
marking to twist, and $N=1$ corresponds to the two boundary-only balanced graphs
of \cite[Example~3.4]{GKT}, which are not critical.)
\end{proof}

\begin{proof}[Proof of $\neg$(i)$\Rightarrow\neg$(iv)]
Assume $\chat_{W}\geq1$, i.e.\ $\frac2r+\frac2s\leq1$; in particular
$r,s\geq3$ and not both equal to $3$, so $r,s\geq3$ with $\max(r,s)\geq4$, and
we may take all interior twists maximal, $(a_{i},b_{i})=(r-2,s-2)$. With $n$ such
markings,
\begin{equation}
\label{eq:maxN}
N=\Big\lfloor\tfrac{n(r-2)}{r}\Big\rfloor+\Big\lfloor\tfrac{n(s-2)}{s}\Big\rfloor-n+1
=\Big(n-\big\lceil\tfrac{2n}{r}\big\rceil\Big)
+\Big(n-\big\lceil\tfrac{2n}{s}\big\rceil\Big)-n+1
=n-\big\lceil\tfrac{2n}{r}\big\rceil-\big\lceil\tfrac{2n}{s}\big\rceil+1,
\end{equation}
using $\lfloor n(r-2)/r\rfloor=\lfloor n-2n/r\rfloor=n-\lceil2n/r\rceil$. We
exhibit $n$ with $N\geq1$, i.e.\ $\lceil2n/r\rceil+\lceil2n/s\rceil\leq n$.

If $\frac2r+\frac2s<1$ (the hyperbolic regime $\chat_{W}>1$), then
$\lceil2n/r\rceil+\lceil2n/s\rceil\leq\frac{2n}{r}+\frac{2n}{s}+2
=n\bigl(\frac2r+\frac2s\bigr)+2$, which is $\leq n$ once
$n\bigl(1-\frac2r-\frac2s\bigr)\geq2$; such $n$ exists since
$1-\frac2r-\frac2s>0$.

If $\frac2r+\frac2s=1$ (the parabolic regime $\chat_{W}=1$: the pairs $(4,4)$,
$(3,6)$, $(6,3)$), choose $n$ with $r\mid2n$ and $s\mid2n$, e.g.\
$n=\operatorname{lcm}(r,s)/2$, an integer for all three parabolic pairs; then
$\lceil2n/r\rceil+\lceil2n/s\rceil=\frac{2n}{r}+\frac{2n}{s}=n$,
so $N=1$.

In either case a primary critical graph exists once $I$ contains $n$ interior
labels of maximal twist, which is possible because each twist is realized
infinitely often in the marking data \cite[\S2.4]{GKT}. By
Proposition~\ref{prop:reduction}(b) the primary invariants over such an $I$
depend on the boundary conditions. This proves $\neg$(iv), and with it
Theorem~\ref{thm:main}.
\end{proof}

\begin{remark}[On Remark~0.9 of \cite{GKT}]
\label{rem:erratum}
\cite[Remark~0.9]{GKT} (arXiv versions~2 and~3) states that the pairs
$(r,s)\in\{(n,2),(2,n),(3,3),(3,4),(4,3)\}$ give simple singularities
``$A_{n-1}$, $E_{6}$, and $E_{8}$''. Two adjustments reconcile the list with the
naming. First, $(3,3)$ appears in the list but $D_{4}$ is not named; $x^{3}+y^{3}$
is the $D_{4}$ singularity (three distinct lines, Milnor number $4$). Second,
$E_{8}$ is named but its pairs $(3,5),(5,3)$ are absent from the list; these are
$x^{3}+y^{5}$, Milnor number $8$. The list, completed and matched to its
names, is
\[
\{(n,2),(2,n):n\geq2\}\cup\{(3,3),(3,4),(4,3),(3,5),(5,3)\},
\]
of types $A_{n-1},D_{4},E_{6},E_{8}$, exactly as in Theorem~\ref{thm:main}(iii).
The theorem shows this list is not merely a set of examples but the complete
answer to the prediction.
\end{remark}

\section{The first wall and the descendent wall}
\label{sec:firstwall}

\begin{proof}[Proof of Theorem~\ref{thm:firstwall}]
The simple case $n_{*}=\infty$ is Theorem~\ref{thm:main}. Assume $\chat_{W}\geq1$
and use \eqref{eq:maxN}, which gives the maximal $N$ over size-$n$ primary
multisets, since $N$ is nondecreasing in the twists and maximal at
$(r-2,s-2)$. Thus $n_{*}(r,s)$ is the least $n$ with
$\lceil2n/r\rceil+\lceil2n/s\rceil\leq n$.

\emph{$\min(r,s)\geq4$.} For $n=1$, $\lceil2/r\rceil+\lceil2/s\rceil=1+1=2>1$
(each $2/r,2/s\leq\frac12$). For $n=2$ and $r,s\geq4$,
$\lceil4/r\rceil+\lceil4/s\rceil=1+1=2=n$, so $N=1$ and $n_{*}=2$.

\emph{$\min(r,s)=3$, say $r=3$, with $\max=s\geq6$.} Here $a_{i}\leq1$, so
$\ell_{1}=\lfloor n/3\rfloor$. For $n=2$,
$\lceil4/3\rceil+\lceil4/s\rceil=2+1=3>2$. For $n=3$,
$\lceil6/3\rceil+\lceil6/s\rceil=2+1=3=n$ (using $s\geq6$), so $N=1$ and
$n_{*}=3$. For $3\le s\le5$ the pair is simple and $n_{*}=\infty$, already
covered.

Finally $\min(r,s)=2$ forces the $A$-type simple case, $n_{*}=\infty$. This
exhausts all $(r,s)$. For $x^{4}+y^{4}$, the witnessing multiset is two
insertions of twist $(r-2,s-2)=(2,2)$, giving $N=1$: the example of
\cite[Examples~4.9 and~5.5]{GKTsurvey}.
\end{proof}

\begin{proof}[Proof of Proposition~\ref{prop:descendent}]
Take $|J|=1$ with any twist $(a,b)$ in the Neveu--Schwarz range. Then
$\ell_{1}=\ell_{2}=0$ and, by \eqref{eq:N},
$N(J,\mathbf d)=0+0-1+1+d=d$, where $d$ is the descendent degree of the single
insertion. Hence $d=1$ gives $N=1$: exactly one critical graph, for every
$(r,s)$. Proposition~\ref{prop:reduction}(b) then gives two systems of invariants
that differ in a descendent invariant. Hence the descendent theory
wall-crosses in every case.  When $(a,b)=(0,0)$, and in particular for
$(r,s)=(2,2)$, where this is the only twist, the critical graph has weight zero
and its generator is a multiple of the torus vector field, so the conclusion
rests on Proposition~\ref{prop:enlargedaction}.
\end{proof}

\begin{remark}
The trichotomy of Theorem~\ref{thm:firstwall} shows the $x^{4}+y^{4}$ example of
\cite{GKTsurvey} is not isolated: two interior insertions already force a wall
for every $(r,s)$ with $r,s\geq4$, and three suffice whenever
$\min(r,s)=3\leq\frac12\max(r,s)$. Combined with
Proposition~\ref{prop:descendent}, the qualitative picture is complete.
\end{remark}

\section{The wall generators are Hamiltonian}
\label{sec:hamiltonian}

Throughout this section $\omega=dx\wedge dy$, and for a polynomial $H$ the
Hamiltonian vector field is
$X_{H}:=(\partial_{y}H)\,\partial_{x}-(\partial_{x}H)\,\partial_{y}$, so that
$\iota_{X_{H}}\omega=dH$.  For $k=(k_1,k_2)\in\ZZ^2_{\ge0}$ put
\begin{equation}
\label{eq:Xab}
X_{k}=X_{k_1,k_2}=x^{k_1}y^{k_2}\bigl((k_2+1)x\partial_{x}-(k_1+1)y\partial_{y}\bigr)
=(k_2+1)x^{k_1+1}y^{k_2}\,\partial_{x}-(k_1+1)x^{k_1}y^{k_2+1}\,\partial_{y},
\end{equation}
and $H_k=x^{k_1+1}y^{k_2+1}$; for a critical graph $\Lambda$ of bidegree $k$ write
$H_\Lambda=H_k$.  The generator of a critical graph of bidegree
$k$ is a multiple of $X_k$, and the exponent $k+(1,1)$ of $H_k$ records the
numbers of its boundary markings.

\begin{lemma}
\label{lem:hamvf}
For every $k$, $X_{k}=X_{H_{k}}$.  In particular $\operatorname{div}X_{k}=0$,
so $X_{k}$ preserves $\omega$, and $X_k$ is tangent to both coordinate axes.
\end{lemma}

\begin{proof}
$\partial_{y}H_{k}=(k_2+1)x^{k_1+1}y^{k_2}$ and
$\partial_{x}H_{k}=(k_1+1)x^{k_1}y^{k_2+1}$, which gives \eqref{eq:Xab}.  A
Hamiltonian vector field is divergence-free, which recovers the
volume-preservation of \cite[p.~89]{GKT}.  The coefficient of $\partial_x$
vanishes on $x=0$ and that of $\partial_y$ on $y=0$.
\end{proof}

\begin{proposition}[Hamiltonian wall generators]
\label{prop:ham}
For $k,l\in\ZZ^2_{\ge0}$,
\begin{equation}
\label{eq:bracketid}
[X_{k},X_{l}]=\bigl((k_2+1)(l_1+1)-(k_1+1)(l_2+1)\bigr)\,X_{k+l},
\end{equation}
equivalently $\{H_k,H_l\}=\bigl((k_2+1)(l_1+1)-(k_1+1)(l_2+1)\bigr)H_{k+l}$
for the Poisson bracket $\{H,K\}:=X_H(K)$ of $\omega$.  Thus $H_k\mapsto X_k$
is an isomorphism of graded Lie algebras from the Poisson algebra of $\omega$
on the span of the monomials $x^py^q$, $p,q\ge1$, onto its image.  After
extension of scalars to $A_{I,\mathrm{sym}}$ it maps the span of the products
$u_{J,\mathbf d}H_\Lambda$, over critical graphs $\Lambda$ of positive weight
with data $(J,\mathbf d)$, onto $\gfrak^{r,s}_{A_{I,\mathrm{sym}}}$, and, with
the products $u_{J,\mathbf d}\,xy$ of the weight-zero critical graphs
adjoined, onto $\widetilde{\gfrak}^{\,r,s}_{A_{I,\mathrm{sym}}}$.  The
$(s,r)$-weighted degree of $H_\Lambda$ is $w(J,\mathbf d)+r+s$.
\end{proposition}

\begin{proof}
Write $X_{k}=f\partial_{x}+g\partial_{y}$ and $X_{l}=F\partial_{x}+G\partial_{y}$
as in \eqref{eq:Xab}.  The $\partial_x$-coefficient of $[X_k,X_l]$ is
\[
X_{k}F-X_{l}f
=(k_2+l_2+1)\bigl((k_2+1)(l_1+1)-(k_1+1)(l_2+1)\bigr)\,x^{k_1+l_1+1}y^{k_2+l_2},
\]
the stated multiple of the $\partial_x$-coefficient of $X_{k+l}$, and the
$\partial_y$-coefficients match in the same way.  The Poisson form follows
from Lemma~\ref{lem:hamvf}, since $X_{k}(H_l)$ is the stated multiple of
$H_{k+l}$.  The generators of $\gfrak^{r,s}$ and of its enlargement are the
products $u_{J,\mathbf d}X_{k}$ over the critical graphs of bidegree $k$
\cite[Definition~4.23]{GKT}, \eqref{eq:enlargedgroup}.  For the grading, the
critical condition \eqref{eq:criticalweight-intro} gives
$sk_1+rk_2=w(J,\mathbf d)$; add $r+s$.
\end{proof}

The coefficient in~\eqref{eq:bracketid} is, up to sign, the determinant of
the boundary exponents $k+(1,1)$ and $l+(1,1)$, as in the tropical vertex Lie
algebra of Kontsevich--Soibelman \cite{KS06,Gross10}.

\begin{remark}[Relation to the tropical vertex algebra]
\label{rem:tropical}
The tropical vertex Lie algebra is the Hamiltonian algebra of the
\emph{logarithmic} form $d\log x\wedge d\log y$, for which
$\{x^{u},x^{v}\}=\langle u,v\rangle\,x^{u+v}$ with $\langle u,v\rangle$ the
determinant, so that Hamiltonians add exponents. Proposition~\ref{prop:ham} uses
instead the \emph{affine-plane} form $dx\wedge dy=xy\,(d\log x\wedge d\log y)$,
under which $\{x^{u},x^{v}\}=\langle u,v\rangle\,x^{u+v-(1,1)}$, and the
Hamiltonian vector fields of the two forms are related by
$X^{\mathrm{aff}}_H=(xy)^{-1}X^{\mathrm{tor}}_H$.  The structure constants
agree, but the exponent of a bracket differs by $(1,1)$, and the two Lie
algebras are not isomorphic as graded Lie algebras.  For instance, in
$\gfrak^{4,6}$ the generators $t_{2,3,0}t_{2,4,0}X_{0,1}$ and
$t_{2,4,0}^2X_{0,2}$ of Table~\ref{tab:census} have nonzero bracket once $I$ carries one label of
twist $(2,3)$ and three of twist $(2,4)$, whereas a graded map sends
$X_{0,1}$ and $X_{0,2}$ to tropical vertex elements with parallel exponents,
which commute.  The analogy that \cite[Remarks~0.8 and~4.20]{GKT} draw
``in spirit'' with the Kontsevich--Soibelman group is therefore an analogy of
structure constants.  The companion paper \cite{companion}
groups the corresponding automorphisms by the rays spanned by
$(k_1+1,k_2+1)$, orders those ray factors by slope, and uses their products to
transport the open potentials.  No result of this paper depends on that
construction.
The tropical open FJRW theory called for in \cite[\S6]{GKTsurvey} is posed
there through tropical moduli, a different route.
\end{remark}

\section{The first wall at the parabolic boundary}
\label{sec:parabolic}

The threshold $\chat_{W}=1$ is not only where wall-crossing begins across the
family; the walls it produces have a uniform shape. Recall that the first
wall of $(r,s)$ is the set of critical graphs over interior multisets of the
minimal size $n_{*}(r,s)$. Combining the count
\eqref{eq:critcount} with the Hamiltonian dictionary of \S\ref{sec:hamiltonian}
determines the first wall of any given $(r,s)$ by a finite enumeration,
recorded in Table~\ref{tab:census} for the parabolic and the smallest
hyperbolic pairs; at the parabolic boundary it collapses to a single wall.

\begin{proposition}[Parabolic first wall]
\label{prop:parabolic}
For any interior multiset $J$ with $N(J,\mathbf 0)=1$, the bidegree of its unique
critical graph is $(k_{1},k_{2})=(r(J),s(J))$, of Euler weight
$s\,r(J)+r\,s(J)$. For the
three parabolic pairs, $(r,s)\in\{(4,4),(3,6),(6,3)\}$ with $\chat_{W}=1$, the
first wall consists of exactly one critical graph, its multiset forced to be
the maximal-twist one: two copies of $(2,2)$ for $(4,4)$, three copies of
$(1,4)$ for $(3,6)$, three copies of $(4,1)$ for $(6,3)$. Its residues vanish,
so its bidegree is $(0,0)$, its weight is zero, and its generator is a multiple
of the torus vector field $x\partial_{x}-y\partial_{y}$ of
\eqref{eq:enlargedgroup}, with Hamiltonian $xy$.  More generally, for the
parabolic pairs every primary critical graph has weight zero, so positive-weight
walls occur exactly in the hyperbolic regime.
\end{proposition}

\begin{proof}
If $N(J,\mathbf 0)=1$ there is a unique critical graph and, by
\cite[Notation~2.40, eqn.~(2.31)]{GKT} read in the indexing of
\eqref{eq:generator}, its bidegree is $(k_{1},k_{2})=(r(J),s(J))$; the weight
statement is the critical condition \eqref{eq:criticalweight-intro}. For each parabolic pair
the minimal critical multiset is uniquely the maximal-twist one: at $n=n_{*}$
the inequality $\ell_{1}+\ell_{2}\geq n$ is attained only with equality and
only at $(a_{i},b_{i})=(r-2,s-2)$ (for $(4,4)$, $\ell_{1},\ell_{2}\leq1$ force
$a_{i}=b_{i}=2$; for $(3,6)$, $\ell_{1}\leq1,\ell_{2}\leq2$ force
$a_{i}=1,b_{i}=4$). There the residues vanish, so $(k_{1},k_{2})=(0,0)$: the
torus vector field, with Hamiltonian $H_{0,0}=xy$ (Lemma~\ref{lem:hamvf}).
For the last statement, Lemma~\ref{lem:bound} gives $N(J,\mathbf 0)\le1$ in
the parabolic regime, and $N(J,\mathbf 0)=1$ forces equality there, hence no
loss in the floors and maximal twists $(r-2,s-2)$; then $r\mid\sum a_i$ and
$s\mid\sum b_i$, so $r(J)=s(J)=0$ and the bidegree is $(0,0)$.  In the hyperbolic
regime, $n$ copies of the maximal twist with $N\ge1$ exist by the proof of
Theorem~\ref{thm:main}, and by \eqref{eq:massweight} every critical graph over
them has Euler weight $n(rs-2r-2s)>0$.
\end{proof}

For the hyperbolic pairs of Table~\ref{tab:census}, the enumeration of the
interior multisets of size $n_*$ with $N\ge1$ shows that the first wall
comprises several critical graphs, some of positive weight and some of
weight zero.

For $(4,4)$ and $(3,6)$ the minimal multisets match the integrality condition
on the exponent of the marginal parameter in Maher's description of the
elliptic primary wall-crossing group \cite[Proposition~5.2.10 and
Example~5.2.13]{Maher}, obtained there for open Saito potentials.

\begin{table}[ht]
\centering
\begin{tabular}{ccccccl}
\toprule
$(r,s)$ & $\chat_{W}$ & regime & $n_{*}$ & graphs & weight zero & wall Hamiltonians\\
\midrule
$(4,4)$ & $1$ & parabolic & $2$ & $1$ & $1$ & --- ($xy$ only)\\
$(3,6)$ & $1$ & parabolic & $3$ & $1$ & $1$ & --- ($xy$ only)\\
$(6,3)$ & $1$ & parabolic & $3$ & $1$ & $1$ & --- ($xy$ only)\\
$(4,5)$ & $11/10$ & hyperbolic & $2$ & $2$ & $1$ & $xy^{2}$\\
$(3,7)$ & $22/21$ & hyperbolic & $3$ & $2$ & $1$ & $xy^{2}$\\
$(4,6)$ & $7/6$ & hyperbolic & $2$ & $4$ & $2$ & $xy^{2},\,xy^{3}$\\
$(5,5)$ & $6/5$ & hyperbolic & $2$ & $5$ & $2$ & $xy^{2},\,x^{2}y,\,x^{2}y^{2}$\\
\bottomrule
\end{tabular}
\caption{The first wall of $x^{r}+y^{s}$: the number of critical graphs at
$n=n_{*}(r,s)$, how many have weight zero (bidegree $(0,0)$, the torus vector
field with Hamiltonian $xy$), and the Hamiltonians $x^{k_{1}+1}y^{k_{2}+1}$ of
the positive-weight wall generators. At the parabolic boundary the first wall
has weight zero; positive-weight walls appear only in the hyperbolic regime.
Computed exactly; see \S\ref{sec:cert}.}
\label{tab:census}
\end{table}

The parabolic pairs are the two-variable Fermat \emph{simple elliptic}
singularities ($x^{4}+y^{4}=\widetilde E_{7}$, $x^{3}+y^{6}=\widetilde E_{8}$).
Along their marginal deformation Saito's primitive form is $dx\wedge dy$
divided by a period of an elliptic curve; in the normalization of
\cite[Theorem~0.3]{GKT} the form $dx\wedge dy$ is kept and the period enters
through the flat coordinates.  The elliptic wall-crossing group of
\cite[\S5.2.2]{Maher} is generated by the torus vector field, and the
explicit jumping of the
$x^{4}+y^{4}$ invariants in \cite[Example~4.9]{GKTsurvey} is its action.

The first hyperbolic Fermat pair with $r=s$ is $x^{5}+y^{5}$, with boundary
monomials $xy^{2},x^{2}y,x^{2}y^{2}$ in Table~\ref{tab:census}; its wall
functions and slope-ordered boundary-ray factorization are computed in
\cite{companion}.

\section{Finite exact checks}
\label{sec:cert}

The proofs above are uniform in $r$ and $s$; the following finite calculations
serve only as independent reproducibility checks.  Substitution in
\eqref{eq:N} for every $2\le r,s\le40$ reproduces the ADE list and the
trichotomy of Theorem~\ref{thm:firstwall}.  The absence assertion outside this
finite range is supplied by Lemma~\ref{lem:bound}, not by enumeration.  The
same calculation reproduces the first-wall data printed in
Table~\ref{tab:census}, including the weight-zero/positive-weight distinction,
and the bidegrees of the single-descendent critical graphs of
Proposition~\ref{prop:descendent}.
Finally, direct polynomial differentiation verifies Lemma~\ref{lem:hamvf}
for $0\le k_1,k_2\le4$ and identity~\eqref{eq:bracketid}, componentwise, for
symbolic $k$ and $l$.\footnote{The exact-arithmetic files are
\texttt{src/census.py} and \texttt{src/gkt\_algebra.py} in the repository cited
in \cite{companion}.}
\section*{Concluding remarks and open directions}
\addcontentsline{toc}{section}{Concluding remarks and open directions}

The companion paper \cite{companion} continues the Hamiltonian identification
by grouping the automorphisms on rays of boundary degree and factoring them in
slope order.  Two further directions remain open for the threshold problem
itself.

\begin{enumerate}
\item \emph{Explicit ADE potentials.} On the wall-crossing-free locus the
primary potential is unambiguous, and by \cite[Theorem~0.3]{GKT} it is the
unfolding of $W$ in flat coordinates for the primitive form $dx\wedge dy$.
The Fermat $D_4$ case is computed in \cite[Example~4.2.3]{Maher}; writing out
the $E_6$ and $E_8$ cases from the classical flat coordinates
\cite{Saito83,NoumiYamada98} and relating
the open invariants to Coxeter and root-system data remain to be done.
\item \emph{The parabolic boundary.} Proposition~\ref{prop:parabolic} pins the
first wall of the parabolic pairs to the torus vector field. What remains is
the analytic side: these are the simple elliptic singularities
$\widetilde E_{7},\widetilde E_{8}$, whose flat structure along the marginal
deformation is governed by modular forms rather than polynomials
\cite[\S4.2.2]{Maher}.  Maher studies the action of the torus vector field on
the modular potential on the B-model side \cite[Example~5.2.12]{Maher};
identifying the corresponding action on the open FJRW invariants, beyond the
$x^{4}+y^{4}$ jumping of \cite[Example~4.9]{GKTsurvey}, is the first
transcendental case of the theory.
\end{enumerate}

\begin{thebibliography}{GKT26}

\bibitem[Arn72]{Arnold72} V.~I.~Arnold,
\emph{Normal forms for functions near degenerate critical points, the Weyl
groups of $A_k$, $D_k$, $E_k$ and Lagrangian singularities},
Funct.\ Anal.\ Appl.\ 6 (1972), no.~4, 254--272.

\bibitem[BB21]{BB21} A.~Basalaev, A.~Buryak,
\emph{Open Saito theory for $A$ and $D$ singularities},
Int.\ Math.\ Res.\ Not.\ 2021, no.~7, 5460--5491. arXiv:1909.00598.

\bibitem[BCT22]{BCT18} A.~Buryak, E.~Clader, R.~J.~Tessler,
\emph{Open $r$-spin theory I: Foundations},
Int.\ Math.\ Res.\ Not.\ 2022, no.~14, 10458--10532. arXiv:2003.01082.

\bibitem[BCT24]{BCT21} A.~Buryak, E.~Clader, R.~J.~Tessler,
\emph{Open $r$-spin theory II: The analogue of Witten's conjecture for $r$-spin
disks}, J.\ Differential Geom.\ 128 (2024), no.~1, 1--75. arXiv:1809.02536.

\bibitem[FJR13]{FJR13} H.~Fan, T.~Jarvis, Y.~Ruan,
\emph{The Witten equation, mirror symmetry, and quantum singularity theory},
Ann.\ of Math.\ 178 (2013), no.~1, 1--106. arXiv:0712.4021.

\bibitem[GKT22]{GKT} M.~Gross, T.~L.~Kelly, R.~J.~Tessler,
\emph{Open FJRW theory and mirror symmetry},
Geom.~Topol.~30 (2026), no.~8, 2779--2960; doi:10.2140/gt.2026.30.2779;
arXiv:2203.02435v3 (2022; revised 2026).

\bibitem[GKT26]{GKTsurvey} M.~Gross, T.~L.~Kelly, R.~J.~Tessler,
\emph{Open enumerative geometries for Landau--Ginzburg models},
arXiv:2602.12707v1 (2026).

\bibitem[GPS10]{GPS10} M.~Gross, R.~Pandharipande, B.~Siebert,
\emph{The tropical vertex}, Duke Math.\ J.\ 153 (2010), no.~2, 297--362.
arXiv:0902.0779.

\bibitem[Gro10]{Gross10} M.~Gross,
\emph{Mirror symmetry for $\mathbb{P}^{2}$ and tropical geometry},
Adv.\ Math.\ 224 (2010), no.~1, 169--245. arXiv:0903.1378.

\bibitem[KS06]{KS06} M.~Kontsevich, Y.~Soibelman,
\emph{Affine structures and non-Archimedean analytic spaces},
in \emph{The unity of mathematics}, Progr.\ Math.\ 244, Birkh\"auser (2006),
321--385. arXiv:math/0406564.

\bibitem[Mah24]{Maher} R.~Maher,
\emph{Predictions in open Fan--Jarvis--Ruan--Witten theory via mirror symmetry,
modularity and wall-crossing}, Ph.D.\ thesis, University of Birmingham, 2024.
\texttt{etheses.bham.ac.uk/id/eprint/15556}.

\bibitem[NY98]{NoumiYamada98} M.~Noumi, Y.~Yamada,
\emph{Notes on the flat structures associated with simple and simply elliptic
singularities}, in \emph{Integrable systems and algebraic geometry}, World Sci.\
(1998), 373--383.

\bibitem[Sai74]{Saito74} K.~Saito,
\emph{Einfach-elliptische Singularit\"aten},
Invent.\ Math.\ 23 (1974), 289--325.

\bibitem[Sai83]{Saito83} K.~Saito,
\emph{Period mapping associated to a primitive form},
Publ.\ Res.\ Inst.\ Math.\ Sci.\ 19 (1983), no.~3, 1231--1264.

\bibitem[Wue26]{companion} B.~J.~Wuebben,
\emph{Boundary-ray factorization and geometric wall-crossing in open FJRW theory
of $x^{r}+y^{s}$},
paper and accompanying computations,
\url{https://github.com/bwuebben/open-fjrw-scattering} (2026).

\end{thebibliography}

\bigskip
\noindent{\sc Bernd J. Wuebben, New York, NY,} \texttt{wuebben@gmail.com}

\end{document}
